\documentclass[letterpaper]{article}
\usepackage[letterpaper,margin=1in]{geometry}
\input{wsp/testing/test-case-library.tex}
\begin{document}
\def\TCID{TC-0070}
\def\TCTitle{Context isolation, fake runtime, and concurrency}
\def\TCPurpose{Verify the M2-owned context portion of WSH-REQ-0070.}
\def\TCPriority{\TCPriorityCritical}
\def\TCRequirementRef{WSH-REQ-0070}
\def\TCDesignRef{ADR-0002, ADR-0005, and M2 runtime design}
\def\TCFeatureRef{Context ownership and deterministic fake runtime}
\def\TCTestTechnique{Concurrency, state isolation, and mock-object testing.}
\def\TCPreconditions{The Windows test host can create two worker threads.}
\def\TCEnvironment{Native Windows x86 or x64 Debug/Release host.}
\def\TCAssumptions{Each context is used by only one thread at a time.}
\def\TCInputData{Same variable name with different values and one scripted
runtime operation.}
\def\TCInitialState{Two independent contexts and one expectation queue.}
\def\TCProcedure{\begin{enumerate}
\item Verify byte-exact scripted output, status, ordering, and mismatch.
\item Set different state in two contexts and inspect isolation.
\item Run independent context update loops on two concurrent threads.
\end{enumerate}}
\def\TCExpectedResult{The fake performs no external effect and consumes
expectations deterministically; contexts do not share state; both threads
complete successfully.}
\def\TCPostConditions{Threads join and every context/runtime owner is released.}
\TCTemplate
\end{document}
